\documentclass[10pt,a4paper]{amsart}
\usepackage[margin=19mm]{geometry}
\usepackage{amsmath,amssymb,mathtools}
\usepackage[scr=rsfs,cal=euler]{mathalfa}
\usepackage{xfrac}
\usepackage[unicode,hidelinks,bookmarks=false]{hyperref}
\newcommand{\ie}{i.e.\ }
\newcommand{\cf}{cf.\ }
\newcommand{\resp}{respectively\ }
\newcommand{\Subsection}{\subsection}
\providecommand{\nobreakdash}{\nobreak\hbox{-}}
\setlength{\emergencystretch}{2em}
\multlinegap0pt
\usepackage{bm}
\usepackage{bbm}
\usepackage{dsfont}
\usepackage{xspace}
\usepackage{cancel}
\usepackage{mathtools}
\usepackage{amsthm}
\makeatletter
\@ifundefined{theorem}{\newtheorem{theorem}{Theorem}}{}
\makeatother
\title[Three Generalizations of Erd\H{o}s Szekeres: $k$-Modal Subse]{Three Generalizations of Erd\H{o}s Szekeres: $k$-Modal Subsequences}
\author{Charles Gong}
\date{Source: arXiv:2508.20360v2 (CC BY 4.0)}
\begin{document}
\maketitle

\noindent\emph{Source and licence.} Three Generalizations of Erd\H{o}s Szekeres: $k$-Modal Subsequences. Version arXiv:2508.20360v2 (CC BY 4.0), distributed under the Creative Commons Attribution 4.0 International licence, as recorded in the arXiv licence metadata for this exact version and shown on the abstract page https://arxiv.org/abs/2508.20360v2. Full rights notice: licenses/arxiv-2508.20360-CC-BY-4.0.txt.\par\smallskip

\noindent\textit{Editorial scope.} Counted unit: the Theorem whose source label is \texttt{None} in the article (source0.tex lines 142--144, source label \texttt{None}), delivered together with the article's own complete original proof of it (source0.tex lines 146--194, one \texttt{proof} environment of 49 lines). No mathematics is written, repaired or re-derived: statement and proof are the article's own text line for line. Required context retained uncounted: none. Every definition, display and label of those lines is retained unchanged. Omitted from this package and not redistributed: 1--141, 145--145, 195--269. Nothing inside a retained excerpt is omitted or altered: every retained body is delivered exactly as the article's lines are written, so no omission receipt of any kind accompanies this package. Citations: the retained excerpts carry no citation command at all, so no bibliography entry of the article is transcribed and no reference is redistributed. Reference adaptation: none was needed and none was made; every reference inside the delivered unit resolves inside it, so no unresolved reference or citation is produced. Printed slips kept literal: no printed word, symbol, hypothesis or number of the counted statement or of its proof was altered. Layout and class: the article's own class is \texttt{article} with packages including graphicx, amsmath, amsfonts, amssymb, babel, amsthm, tikz, thm-restate, geometry, hyperref; the wrapper is the lane's pinned \texttt{amsart} editorial wrapper at 10pt on A4, which restates the theorem-like environments the retained text needs and adds the article's own macro definitions that the retained text needs (none), with extra wrapper packages (bm, bbm, dsfont, xspace, cancel, mathtools). The delivered numbering is the unit's own. Assets: no retained span contains a figure environment, a graphics-inclusion command, a PDF-page inclusion, a table or a TikZ picture; no image file is copied into this package, the package declares no asset dependency and no image file is redistributed with it. Licence: version arXiv:2508.20360v2 of this article is distributed under the Creative Commons Attribution 4.0 International licence, as recorded in the arXiv licence metadata for this exact version and shown on the abstract page https://arxiv.org/abs/2508.20360v2. A full rights notice is delivered at licenses/arxiv-2508.20360-CC-BY-4.0.txt. Characters: every retained excerpt is pure ASCII, so the delivered engine, XeLaTeX with its default Latin Modern text font, reports no missing character.
\begin{theorem}
In any length $n$ permutation $(a_1, a_2, \ldots, a_n)$, there exists an $a_i$ such that there exists both an increasing and a decreasing first $k$-modal subsequence of lengths at least $\sqrt{2kn}$ (up to some constants) ending at $a_i$. 
\end{theorem}

\begin{proof}
We again proceed by induction. $k = 0$ is trivial. We start at $k = 1$ for illustrative purposes. \\

Define $x(a_i)$ to be the maximum length of an increasing subsequence ending at $a_i$. Define $y(a_i)$ to be the maximum length of a decreasing subseqeuence ending at $a_i$. Note this definition of $y$ differs than the one given in Theorem 1. However, using a similar argument in Theorem 1 it can be shown that the function $f : [n] \rightarrow [n] \times [n]$ defined by $f(a_i) = (x(a_i), y(a_i))$ is injective. Define $N$ to be the maximum number such that there exists a point $a_i$ with an increasing and a decreasing $1$-modal subsequence both of length at least $N$ ending at $a_i$. Note that $0$-modal subsequences, i.e. increasing or decreasing subsequences, also count as $1$-modal subsequences. In fact, any $0$-modal subsequence is both an increasing and a decreasing first $1$-modal subsequence. So in particular $x(a_i), y(a_i) \leq N$ as well (an increasing subsequence of length $x(a_i)$ is both an increasing and a decreasing first $1$-modal subsequence of length $x(a_i)$ ending at $a_i$, so $x(a_i) \leq N$ and similarly for $y(a_i)$), and the picture to have in mind is that $f$ is mapping the points $a_i$ to a square box $\{(x,y) \in \mathbb{N}^2 : 1 \leq x, y \leq N \}$. \\

We note a key fact. Note that if there exists $(a,b) \in [N]^2$ such that 

$$|\{a_i : x(a_i) = a \text{ and } y(a_i) \leq b\}| > N+1-a$$

and 

$$|\{a_i : y(a_i) = b \text{ and } x(a_i) \leq a\}| > N+1-b$$

then there exists a point $a_i$ such that there exists both an increasing and a decreasing first $1$-modal subsequence of lengths at least $N+1$ ending at $a_i$, which would contradict the definition of $N$. We call the inequalities above \emph{the condition}. So if the condition is satisfied (I will use "triggered" interchangeably with "satisfied") for some $(a,b) \in [N]^2$, there exists a point $a_i$ such that there exists both an increasing and a decreasing first $1$-modal subsequence of lengths at least $N+1$ ending at $a_i$. \\

To see this, define $A_{(a,b)} = \{a_i : x(a_i) = a \text{ and } y(a_i) \leq b\}$ and $B_{(a,b)} = \{a_i : y(a_i) = b \text{ and } x(a_i) \leq a\}$. Now order the elements $a_i \in A_{(a,b)}$ according to $y(a_i)$, so we have $y(a_{i_1}) < y(a_{i_2}) < \ldots < y(a_{i_m})$ where $m > N+1-a$ (we know the $y$-values are distinct because they all share the same $x$-value, and $f$ is injective). Note that $(a_{i_1}, a_{i_2}, \ldots, a_{i_m})$ forms a decreasing subsequence because they share the same $x$-values. Also note by assumption $x(a_{i_1}) = a$. Thus, we know that there's an increasing first $1$-modal subsequence of length $a + m - 1 \geq a + (N+1-a) = N+1$ ending at $a_{i_m}$. Using a similar argument and ordering the elements in $B_{(a,b)}$ as $(a_{j_1}, a_{j_2}, \ldots, a_{j_p})$, there's a decreasing first $1$-modal subsequence of length $\geq N+1$ ending at $a_{j_p}$. \\

Case 1: if $i_m \leq j_p$ then we must have $a_{i_m} \geq a_{j_p}$. Otherwise assume for the sake of contradiction that $a_{i_m} < a_{j_p}$. Note that this implies $a = x(a_{i_m}) < x(a_{j_p})$ (the first equality is from $a_{i_m} \in A_{(a,b)}$). However, by assumption $x(a_{j_p}) \leq a$ because $a_{j_p} \in B_{(a,b)}$, a contradiction. Thus $a_{i_m} \geq a_{j_p}$. Note that as mentioned above, there's a decreasing first $1$-modal subsequence of length $\geq N+1$ ending at $a_{j_p}$, and also there's an increasing first $1$-modal subsequence of length $\geq N+1$ ending at $a_{i_m}$. Since $i_m \leq j_p$ and $a_{i_m} \geq a_{j_p}$, we can append $a_{j_p}$ to this increasing first $1$ modal subsequence, to get an increasing first $1$ modal subsequence of length $\geq N+1$ ending at $a_{j_p}$. \\

Case 2: if $i_m > j_p$ or $j_p < i_m$, then we must have $a_{j_p} \leq a_{i_m}$. Otherwise assume for the sake of contradiction that $a_{j_p} > a_{i_m}$ Note that this implies that $b = y(a_{j_p}) < y(a_{i_m}) \leq b$, a contradiction. Note that as mentioned above, there's an increasing first $1$-modal subsequence of length $\geq N+1$ ending at $a_{i_m}$, and also there's a decreasing first $1$-modal subsequence of length $\geq N+1$ ending at $a_{j_p}$. Since $j_p < i_m$ and $a_{j_p} \leq a_{i_m}$, we can append $a_{i_m}$ to this decreasing first $1$ modal subsequence, to get a decreasing $1$ modal subsequence of length $\geq N+1$ ending at $a_{i_m}$. \\

To recapitulate, we showed that for any $(a,b) \in [N]^2$, if $|A_{(a,b)}| > N+1-a$ and $|B_{(a,b)}| > N+1-b$, then we have an element $a_i$ which has both an increasing and decreasing first $1$ modal subsequence of length $\geq N + 1$ ending at it. \\

Now we consider the question, how many points can we place in the $N$ by $N$ square such that we avoid triggering the condition above? I.e., what is the maximum number of points we can place in the $N$ by $N$ square such that for all $(a,b) \in [N]^2$, we have either $|A_{(a,b)}| \leq N+1-a$ or $|B_{(a,b)}| \leq N+1-b$? We claim that it is $\frac{N(N+1)}{2}$, but again we approximate this by the area of the triangle defined by the convex hull of $T = \{(x,y) \in [N]^2 : x + y \leq N+1\}$ which is $\frac{N^2}{2}$. \\

We note that the set of points in the triangle $T$ does not trigger the condition, and $|T| = \frac{N(N+1)}{2}$. Now we show any set of points $S \subseteq [N]^2$ that does not trigger the condition has at most $\frac{N(N+1)}{2}$ points. \\

To see this, we claim you can fit all the points into triangle $T$ by shifting them leftward or downward. Take note of the points $(a,b) \in S$ that satisfy $|A_{(a,b)}| > N+1-a$. Note for such a point, in order to avoid triggering the condition at $(a,b)$, we must have $|B_{(a,b)}| \leq N+1-b$. Call the set of such points $C$. Consider the procedure below, which may not be possible: \\

Step 1: For each point $(x,y) \in S$ such that there exists $(a,b) \in C$ such that $x \leq a$ and $y \geq b$, shift $(x,y)$ to the left to fit it inside $T$. \\

Step 2: For all the other points, shift them down into $T$. \\

We claim that this procedure is always possible, and it shifts all the elements of $S$ into $T$. For the points that are shifted leftward in step 1, by assumption there exists $(a,b) \in C$ such that $x \leq a$ and $y \geq b$. By definition of $C$, we know that $|A_{(a,b)}| > N+1-a$. In particular, consider the point $(a, y)$. Since $|A_{(a,y)}| \geq |A_{(a,b)}| > N+1-a$ we must have $|B_{(x,y)}| \leq |B_{(a,y)}| \leq N+1-y$. In particular, this means we can shift everything to the left of $(x,y)$ also including $(x,y)$ itself into $T$. \\

Now consider $(x,y) \in S$ such that there is no $(a,b) \in C$ where $x \leq a$ and $y \geq b$. In particular, this means $(x,y) \not\in C$. Thus $|A_{(x,y)}| \leq N+1-x$, and we can shift all the points below $(x,y)$ including $(x,y)$ itself downward into $T$, but this argument only holds if we didn't shift some points leftward into column $x$ in step 1. But this can only happen if we shift some point $(x',y')$ leftward into column $x$, and there exists $(a,b) \in C$ such that $x' \leq a$ and $y' \geq b$. Note $x \leq x' \leq a$ and $y \geq y' \geq b$, contradicting our original assumption. \\

Thus, approximating the number integer lattice points in $T$ by $\frac{N^2}{2}$, we get that if the image of $f$ on $\{i \in [n] : a_i\}$ does not trigger the condition for any $(a,b)$, we have $n \leq \frac{N^2}{2}$, or $\sqrt{2n} \leq N$. Any number of points greater than $\frac{N^2}{2}$ would trigger the condition and result in both an increasing and a decreasing $1$-modal subsequence of length $\geq N+1$ ending at the same point, contradicting the definition of $N$, the maximum number such that there exists a point $a_i$ with an increasing and a decreasing $1$-modal subsequence both of length at least $N$ ending at $a_i$.  \\

Now we step into the induction step. Assume that in a permutation $(a_1, a_2, \ldots, a_n)$ there exists an $a_i$ such that there exists increasing and decreasing first $k$-modal subsequeunce of lengths at least $\sqrt{2kn}$ ending at $a_i$. We want to show there exists an $a_i$ such that there exists increasing and decreasing first $k+1$-modal subsequeunce of lengths at least $\sqrt{2(k+1)n}$ ending at $a_i$. \\

Define $x(a_i)$ to be the maximum length of an increasing $k$-modal subsequence ending at $a_i$, and $y(a_i)$ to be the maximum length of a decreasing $k$-modal subsequence ending at $a_i$. The function $f : [n] \rightarrow [n] \times [n]$ defined by $f(a_i) = (x(a_i), y(a_i))$ is an injection, which can be proved by a similar argument in theorem 1. Define $N$ as the maximum number such that there exists a point $a_i$ with an increasing and a decreasing $k+1$-modal subsequence both of length at least $N$ ending at $a_i$. Again, note a $k$-modal subsequence is both an increasing and a decreasing first $k+1$-modal subsequence, so $f$ maps the permutation into an $N$ by $N$ square. \\

By the induction hypothesis, there's at most $\frac{kN^2}{2(k+1)^2}$ points $(x(a_i), y(a_i))$ such that $x(a_i) \leq \frac{kN}{k+1}$ or $y(a_i) \leq \frac{kN}{k+1}$. This is because $\sqrt{2k (\frac{kN^2}{2(k+1)^2})} = \frac{kN}{k+1}$. Note that the region $\{(x,y) \in \mathbb{R}^2 : x \leq \frac{kN}{k+1} \text{ or } y \leq \frac{kN}{k+1}\}$ is the complement of a square $R$ in the upper right corner, $\{(x,y) \in \mathbb{R}^2 : x > \frac{kN}{k+1} \text{ and } y > \frac{kN}{k+1}\}$. That square has area $(N - \frac{kN}{k+1})^2 = \frac{N^2}{(k+1)^2}$. Using a similar argument in the $k=1$ case, one can show that if the condition is satisfied for some $(a,b) \in [N]^2$, there exists a point $a_i$ such that there exists both an increasing and a decreasing first $k+1$-modal subsequence of lengths at least $N+1$ ending at $a_i$, contradicting the definition of $N$, the maximum number such that there exists a point $a_i$ with an increasing and a decreasing $k+1$-modal subsequence both of length at least $N$ ending at $a_i$. Note that any subset $S \subseteq R$ satisfies a similar conditional implication, and thus using a similar argument we get that if the condition is not triggered, there is at most $\frac{1}{2}(N - \frac{kN}{k+1})^2 = \frac{N^2}{2(k+1)^2}$ points $(x(a_i), y(a_i))$ in $R$. Thus the number of points in total $n$ satisfies the inequality

$$n \leq \frac{kN^2}{2(k+1)^2} + \frac{1}{2}\frac{N^2}{(k+1)^2} = \frac{N^2}{2(k+1)}$$

and so $\sqrt{2(k+1)n} \leq N$. 
\end{proof} 

\end{document}
